% =============================================================================
% PAPER 3 (Part II) -- Information limits and temporal design for ghost imaging
% through dynamic scattering media.
%
% Companion to Part I (the existing identifiability manuscript). Scaffolded per
% the R26 ruling, docs/ROUND63_GPT_ROUND26_RULING_RAW.md, which is NORMATIVE.
% Spine follows ruling section 24 (eight sections). Provenance of every block
% (verbatim / TODO / open-lemma) is recorded in NOTES.md.
%
% -----------------------------------------------------------------------------
% FENCE -- THE SEVEN CLAIMS THIS PROGRAM MUST NOT MAKE (ruling section 23).
% Future editors: do NOT let any of the following creep into the manuscript.
%   1. No universal image-quality bound from Fisher information alone.
%   2. No claim that a neural network is close to the ceiling without measuring
%      its gap to the bound.
%   3. No universal two-dimensional temporal ridge.
%   4. No claim that temporal randomization is always optimal.
%   5. No claim that the scalar gain model describes arbitrary dynamic turbid
%      media.
%   6. No claim that identifiability thresholds imply statistical efficiency.
%   7. No claim that an information-optimal schedule is automatically
%      PSNR-optimal under nonlinear regularization.
% The clean hierarchy that replaces them (ruling section 23):
%   physical scalar-gain hypothesis -> identifiability -> observed information
%     -> Bayes risk bound -> temporal design.  Each arrow requires assumptions;
%   none may be skipped.
% -----------------------------------------------------------------------------

% Single-column article (not paper2's twocolumn): the O3/O4 matrix displays are
% too wide for a column. All other conventions (macros, amsthm, unsrt bib, \SPH
% placeholders) follow paper2/main_m1.tex. See NOTES.md.
\documentclass[11pt]{article}
\usepackage{amsmath,amssymb,bm}
\usepackage{amsthm}
\usepackage{graphicx}
\usepackage[margin=2.5cm]{geometry}
\usepackage{xcolor}
\usepackage{hyperref}
\usepackage{placeins}

\newcommand{\SPH}[1]{\textcolor{red}{[#1]}}
\newcommand{\TODO}[1]{\textcolor{blue}{[TODO: #1]}}
\newcommand{\Cov}{\operatorname{Cov}}
\newcommand{\diag}{\operatorname{diag}}
\newcommand{\tr}{\operatorname{tr}}
\newcommand{\E}{\mathbb{E}}

\newtheorem{theorem}{Theorem}
\newtheorem{corollary}{Corollary}
\newtheorem{proposition}{Proposition}
\newtheorem{conjecture}{Conjecture}
\newtheorem{hypothesis}{Hypothesis}
\theoremstyle{remark}
\newtheorem{remark}{Remark}

\title{Information limits and temporal design for ghost imaging through
dynamic scattering media}
\author{\SPH{author block --- user decision}}
\date{}

\begin{document}
\maketitle

\begin{abstract}
Ghost imaging through a dynamic scattering medium replaces a fixed forward
operator with a time-varying gain, and correcting for that gain is the central
obstacle to quantitative reconstruction. We derive an exact
missing-information identity showing that a hidden multiplicative gain path
removes precisely the conditional covariance of the complete object score
given the bucket record. In the Poisson bucket channel this loss specializes to
the posterior covariance of a hidden gain-weighted illumination vector, which
yields an estimator-independent Bayes-risk ceiling on any correction --- neural
or otherwise --- under the declared gain model and prior. A Schur-complement
analysis then characterizes the information-optimal temporal schedules:
chopping, interleaving, and randomization help exactly when they make the
object score orthogonal, in the noise metric, to the admitted gain-drift
modes. All statements are conditioned on a stated scalar-gain hypothesis and
its physical validity boundary; a resource-conditioned temporal operating
surface is left as a conditional program rather than a universal claim.
\end{abstract}

% =============================================================================
\section{Introduction}
\label{sec:intro}
Single-pixel and computational ghost imaging (GI) reconstruct a scene from the
responses of one bucket detector to a sequence of structured illumination
patterns~\cite{duarte2008singlepixel,shapiro2008computational,edgar2019principles}.
When the light passes through a dynamic scattering medium, the effective
forward operator acquires a time-varying gain, and a substantial experimental
line corrects for that
gain~\cite{xiao2022temporal,pengchen2023,zhou2023dual,pengchen2024apl,hao2025photonlimited}.
That literature establishes methods and experiments; it does not establish an
exact observed-versus-complete Fisher decomposition, an estimator-independent
Bayes-risk ceiling, or an information-optimal temporal schedule under
correlated gain. Those are the residual questions addressed here.

This manuscript is the information-and-design companion to
Part~I~\cite{companion}, which asks when the object and gain class can be
separated at all. Part~I settles identifiability; the present paper asks, once
separation is possible, how much object information survives hidden gain, what
estimator-independent risk floor follows, and how patterns should be ordered in
time. The two are companion questions, not sections of one paper, and are not
merged. The argument follows a single hierarchy in which each arrow requires
its own assumptions and none may be skipped:
\begin{equation*}
\text{scalar-gain hypothesis}
\rightarrow \text{identifiability}
\rightarrow \text{observed information}
\rightarrow \text{Bayes risk bound}
\rightarrow \text{temporal design}.
\end{equation*}

% =============================================================================
\section{Physical hypothesis MG and channel model}
\label{sec:mg}
The entire program rests on a physical hypothesis: that the dynamic medium acts,
on the illuminated object subspace, as a single scalar multiplier per frame plus
a controlled residual. This section states the hypothesis, its sufficient
physical conditions, and the regimes in which it breaks.

\begin{hypothesis}[MG --- scalar multiplicative gain]
\label{hyp:mg}
Write the dynamic medium's intensity-transfer functional at time $t$ as
$\mathcal K_t$. The scalar multiplicative model is a valid approximation on the
experiment's illuminated object subspace $\mathcal S$ only if there exist a
fixed functional $\mathcal K_0$, a positive scalar $a(t)$, and a small residual
$E_t$ such that
\begin{equation}
\mathcal K_t|_{\mathcal S}
= a(t)\,\mathcal K_0|_{\mathcal S} + E_t,
\qquad
\|E_t z\| \le \varepsilon\, a(t)\,\|\mathcal K_0 z\|
\quad \forall z \in \mathcal S.
\tag{MG}\label{eq:mg}
\end{equation}
After frame integration, $a_n$ is either the quasi-static gain or the properly
weighted frame average.
\end{hypothesis}

Plausible sufficient physical conditions for Hypothesis~\ref{hyp:mg} are:
\begin{enumerate}
\item the bucket collects enough output speckles/modes that temporal medium
variation mainly changes total throughput rather than the normalized spatial
transfer;
\item the object and medium are static enough within a pattern exposure for a
scalar frame average to be meaningful;
\item the pattern does not perturb the medium;
\item incoherent/intensity-linear propagation is an adequate description, or
coherent cross terms average out at the bucket;
\item additive scattered background is separately modeled;
\item the same collection aperture, polarization, and wavelength distribution
apply across patterns.
\end{enumerate}

The scalar model fails when:
\begin{itemize}
\item the dynamic medium changes the spatial transmission matrix across object
coordinates;
\item only a few speckles are collected, producing pattern-dependent
non-Gaussian fluctuations;
\item coherent bucket integration retains interference cross terms that depend
on the displayed pattern;
\item the medium evolves substantially within a frame and cannot be represented
by one scalar average;
\item object motion and medium motion couple to the pattern order;
\item polarization, wavelength, or path-length composition changes with time;
\item multiple scattering creates an additive or convolutive component rather
than common multiplicative throughput.
\end{itemize}

\begin{remark}[Scope of the scalar model]
\label{rem:mg-scope}
In the broken regimes the correct latent object is a vector/matrix-valued
transfer process. The general missing-information identity of
Section~\ref{sec:o1} survives, but the low-dimensional correction theory and
the Poisson hidden-exposure corollary of Section~\ref{sec:o1p} do not. A second
risk is gauge singularity: in a non-identifiable square design both complete and
observed information can be singular. The identity remains true but cannot
manufacture identifiability, so Part~I's thresholds~\cite{companion} must be
assumed before any inverse information or Bayes risk is interpreted.
\end{remark}

% =============================================================================
\section{The exact gain-path missing-information theorem}
\label{sec:o1}
Let $X$ denote the object parameter, $L$ the latent gain path, and $Y$ the
observed bucket record. Assume: (1)~$p(l)$ does not depend on $x$;
(2)~$p(y,l\mid x)=p(l)\,p(y\mid x,l)$ is dominated and differentiable in $x$;
(3)~differentiation may pass under the integral; (4)~the complete score is
square-integrable. Define the complete object score
\begin{equation*}
U_x(Y,L) = \nabla_x \log p(Y,L\mid x) = \nabla_x \log p(Y\mid x,L).
\end{equation*}

\begin{theorem}[O1 --- scattering missing-information identity]
\label{thm:o1}
The observed score is
\begin{equation}
\nabla_x \log p(Y\mid x)
= \E\{ U_x(Y,L) \mid Y, x \}.
\tag{O1.1}\label{eq:o1-obs}
\end{equation}
The complete and observed Fisher matrices satisfy
\begin{equation}
I_Y(x) = I_{Y,L}(x) - \E_Y\!\left[ \Cov\{ U_x(Y,L) \mid Y, x \} \right].
\tag{O1.2}\label{eq:o1-fisher}
\end{equation}
The missing term is positive semidefinite. It vanishes if and only if the
complete object score is measurable from the observed record.
\end{theorem}

\begin{proof}
Differentiate the marginal likelihood:
\begin{equation*}
\nabla_x p(y\mid x)
= \int \nabla_x p(y,l\mid x)\, dl
= \int p(y,l\mid x)\, U_x(y,l)\, dl.
\end{equation*}
Division by $p(y\mid x)$ gives~\eqref{eq:o1-obs}. Scores have mean zero, so the
law of total covariance gives
\begin{equation*}
\Cov(U_x) = \Cov\{ \E(U_x\mid Y) \} + \E\{ \Cov(U_x\mid Y) \}.
\end{equation*}
The first covariance is the observed Fisher matrix and the left side is the
complete Fisher matrix, proving~\eqref{eq:o1-fisher}.
\end{proof}

This proof is classical Louis theory~\cite{louis1982}. The following
programmable-GI specialization is the important part.

% =============================================================================
\section{Poisson hidden-exposure corollary and Gaussian counterpart}
\label{sec:o1p}
\subsection{Exact Poisson bucket corollary}
\label{sec:o1p-poisson}
Let the rows of $M\in\mathbb R_+^{N\times K}$ be $m_n^\top$, let
\begin{equation*}
s = Mx, \qquad s_n > 0, \qquad a_n = e^{l_n} > 0,
\end{equation*}
and conditionally on the entire correlated gain vector $a$, let
$Y_n\mid a,x \sim \operatorname{Poisson}(a_n s_n)$ independently across frames.
The prior on $a$ may have arbitrary temporal correlation and need not be
Gaussian. The complete score is
\begin{equation}
U_x = M^\top D_s^{-1} Y - M^\top a.
\tag{O1.3}\label{eq:o1p-score}
\end{equation}
Conditionally on $a$,
$I_{Y\mid a}(x) = M^\top \diag(a_n/s_n)\, M$.

\begin{corollary}[O1-P --- hidden exposure-vector identity]
\label{cor:o1p}
Under the conditionally Poisson multiplicative-gain model,
\begin{equation}
I_Y(x)
= M^\top \diag\!\left( \frac{\E a_n}{s_n} \right) M
- M^\top \E_Y\{ \Cov(a\mid Y,x) \} M.
\tag{O1.4}\label{eq:o1p-fisher}
\end{equation}
\end{corollary}

\begin{proof}
Average the conditional Fisher matrix for the complete-information term. Given
$Y$, the first term of~\eqref{eq:o1p-score} is fixed, hence
$\Cov(U_x\mid Y,x) = M^\top \Cov(a\mid Y,x)\, M$. Insert this
into~\eqref{eq:o1-fisher}.
\end{proof}

\paragraph{What plays the role of live time.}
For the dead-time channel of Part~I's dead-time line, hidden active time is the
latent exposure to the rate parameter. Here the corresponding object is
\begin{equation}
G_M(a) = M^\top a = \sum_{n=1}^N a_n m_n,
\tag{O1.5}\label{eq:o1p-exposure}
\end{equation}
the \emph{gain-weighted illumination/exposure vector}. The information loss is
exactly the posterior covariance of this hidden exposure vector,
\begin{equation*}
I_{\mathrm{missing}} = \E_Y \Cov\{ G_M(a)\mid Y,x \}.
\end{equation*}
This is the cleanest bridge from Part~I's live-time identity to
dynamic-scattering GI. Three consequences are immediate:
\begin{enumerate}
\item temporal correlation enters through the full posterior covariance of $a$,
not through one scalar variance;
\item pattern ordering matters because it rotates gain uncertainty through
$M^\top(\cdot)M$;
\item a second monitor detector or other gain side information can only reduce
the missing term by conditioning more strongly.
\end{enumerate}

\subsection{Gaussian readout corollary}
\label{sec:o1p-gauss}
For analog bucket data $Y\mid a,x \sim \mathcal N(D_a Mx, R)$ with $R\succ0$, the
complete score and information are
\begin{equation*}
U_x = M^\top D_a R^{-1}(Y - D_a Mx),
\qquad
I_{Y,a}(x) = \E_a[\, M^\top D_a R^{-1} D_a M \,].
\end{equation*}
Identity~\eqref{eq:o1-fisher} remains exact, but the missing term is
\begin{equation*}
\E_Y \Cov\!\left[ M^\top D_a R^{-1}(Y - D_a Mx) \mid Y,x \right],
\end{equation*}
not merely $M^\top \Cov(a\mid Y)\, M$. The unusually clean hidden-exposure
formula~\eqref{eq:o1p-fisher} is therefore a property of the conditionally
Poisson multiplicative channel, not a universal gain law.

% =============================================================================
\section{Universal correction ceiling via van Trees}
\label{sec:o3}
\subsection{Estimator-independent Bayesian bound}
\label{sec:o3-bound}
Let the object have differentiable prior $\pi(x)$, with prior information
$I_\pi = \E_\pi[\nabla\log\pi(X)\,\nabla\log\pi(X)^\top]$. Under the usual van
Trees boundary conditions~\cite{vantrees,tichavsky1998}, every estimator
$\widehat x(Y)$ satisfies
\begin{equation}
\E[(\widehat x - X)(\widehat x - X)^\top]
\succeq \left[ I_\pi + \E_X I_Y(X) \right]^{-1}.
\tag{O3.1}\label{eq:o3-vantrees}
\end{equation}
Combining with Theorem~\ref{thm:o1} gives
\begin{equation}
\E[(\widehat x - X)(\widehat x - X)^\top]
\succeq \left[ I_\pi + \E I_{Y,L}(X) - \E\Cov(U_X\mid Y,X) \right]^{-1}.
\tag{O3.2}\label{eq:o3-combined}
\end{equation}
For a task matrix $W\succeq0$,
\begin{equation}
\E\| W^{1/2}(\widehat x - X) \|_2^2
\ge \tr\!\left\{ W\,[ I_\pi + \E I_Y(X) ]^{-1} \right\}.
\tag{O3.3}\label{eq:o3-task}
\end{equation}
This covers a U-Net, an untrained network, a joint MAP estimator, or any other
estimator that receives the same bucket record and side information.

\begin{remark}[What may be claimed]
\label{rem:o3-may}
Under the declared gain prior and observation model, no estimator using the same
measurements can achieve Bayes MSE below~\eqref{eq:o3-task}.
\end{remark}

\begin{remark}[What may not be claimed]
\label{rem:o3-maynot}
The following are outside the scope of~\eqref{eq:o3-task}:
\begin{itemize}
\item the bound does not show that a particular network is close to optimal;
\item it does not directly bound PSNR, SSIM, visibility, or CNR without an
explicit transformation/loss argument;
\item if a network receives additional calibration data, reference-detector
data, or a mismatched training distribution, those must be included in the
experiment and prior;
\item a lower bound alone does not prove that ``more training is pointless'';
that conclusion requires the achieved risk to be close to the bound.
\end{itemize}
\end{remark}

\subsection{Connection to Part~I's static-correction ceiling}
\label{sec:o3-vbl}
Use the small-log-gain linearization
\begin{equation*}
a \simeq \mathbf 1 + \ell,
\qquad
Y = s + D_s \ell + \epsilon,
\qquad
s = Mx,
\end{equation*}
with $\ell \sim \mathcal N(0, C_\ell)$ and $\epsilon \sim \mathcal N(0, R)$.
Suppose a static correction removes a declared drift subspace with projector
$P$. The residual gain is $\ell_\perp = (I-P)\ell$ and
$C_\perp = (I-P)C_\ell(I-P)^\top$. The marginalized model is
\begin{equation}
Y\mid x \sim \mathcal N(s, \Sigma_x),
\qquad
\Sigma_x = R + D_s C_\perp D_s.
\tag{O3.4}\label{eq:o3-marg}
\end{equation}
Its exact Gaussian Fisher matrix is
\begin{equation}
[I_Y]_{jk}
= (\partial_j s)^\top \Sigma_x^{-1} (\partial_k s)
+ \tfrac12 \tr\!\left( \Sigma_x^{-1} \partial_j\Sigma_x\, \Sigma_x^{-1}
\partial_k\Sigma_x \right).
\tag{O3.5}\label{eq:o3-gaussfisher}
\end{equation}
If $C_\perp = v\bar C$ and $v\to0$, the covariance-derivative term is $O(v^2)$,
while
\begin{equation}
I_Y = M^\top R^{-1} M
- v\, M^\top R^{-1} D_s \bar C D_s R^{-1} M
+ O(v^2).
\tag{O3.6}\label{eq:o3-vexpand}
\end{equation}

\begin{remark}[The $vB_L$ term --- derivation TODO]
\label{rem:o3-todo}
A scalarized risk or relative-MSE expansion of~\eqref{eq:o3-vexpand} yields a
term of the form $v B_L$, so Part~I's static-correction ceiling can be
interpreted as the first-order scalarization of the exact
missing-information/Bayes-risk architecture --- but only under the small-gain
Gaussian model, the declared correction projector, and the chosen task
scalarization. The ruling sketches this connection but does not carry out the
scalarization. \TODO{derive the $vB_L$ term explicitly: fix the task
scalarization $W$, expand $\tr\{W\,[I_\pi+\E I_Y]^{-1}\}$
from~\eqref{eq:o3-vexpand} to first order in $v$, and identify $B_L$ with
Part~I's static-correction coefficient. This is a theorem path, not an
automatic identity between manuscripts.}
\end{remark}

% =============================================================================
\section{Nuisance-orthogonal temporal design}
\label{sec:o4}
\subsection{Linearized drift model}
\label{sec:o4-model}
Let the gain path be represented locally by $\ell = H\beta + r$, where the
columns of $H\in\mathbb R^{N\times p}$ are declared low-frequency or OU/KL drift
modes, $\beta$ has prior precision $\Lambda_\beta$, and $r$ is residual gain
noise. At $\beta=0$, use the local Gaussian model
\begin{equation}
y = M_\pi x + D_{s_\pi} H\beta + \epsilon,
\qquad
s_\pi = M_\pi x,
\qquad
\epsilon \sim \mathcal N(0,R),
\tag{O4.1}\label{eq:o4-model}
\end{equation}
where $\pi$ denotes the temporal order/schedule of the programmable patterns.
The joint information matrix for $(x,\beta)$ is
\begin{equation*}
\mathcal I(\pi) =
\begin{bmatrix}
I_{xx} & I_{x\beta} \\
I_{\beta x} & I_{\beta\beta}
\end{bmatrix},
\end{equation*}
with
\begin{equation*}
I_{xx} = M_\pi^\top R^{-1} M_\pi,
\quad
I_{x\beta} = M_\pi^\top R^{-1} D_{s_\pi} H,
\quad
I_{\beta\beta} = H^\top D_{s_\pi} R^{-1} D_{s_\pi} H + \Lambda_\beta.
\end{equation*}
The efficient/Bayesian information for the object after treating drift as
nuisance is the Schur complement
\begin{equation}
I_{x\mid\mathrm{gain}}(\pi)
= I_{xx}(\pi) - I_{x\beta}(\pi)\, I_{\beta\beta}(\pi)^{-1}\, I_{\beta x}(\pi).
\tag{O4.2}\label{eq:o4-schur}
\end{equation}

\subsection{The nuisance-orthogonal schedule}
\label{sec:o4-theorem}
\begin{theorem}[O4-A --- nuisance-orthogonal schedule]
\label{thm:o4a}
Assume $I_{\beta\beta}\succ0$. For schedules having the same oracle object
information $I_{xx}$,
\begin{equation*}
I_{x\mid\mathrm{gain}} \preceq I_{xx},
\end{equation*}
with equality if and only if
\begin{equation}
M_\pi^\top R^{-1} D_{M_\pi x} H = 0.
\tag{O4.3}\label{eq:o4-orth}
\end{equation}
\end{theorem}

\begin{proof}
The subtracted term in~\eqref{eq:o4-schur} is of the form $AA^\top$ with
$A = I_{x\beta} I_{\beta\beta}^{-1/2}$, and is therefore positive semidefinite.
It vanishes exactly when $I_{x\beta}=0$, proving the result.
\end{proof}

\paragraph{Exact moment conditions.}
The optimal temporal schedule does not merely make gains ``stationary.'' It
makes the object score orthogonal, in the noise metric, to every admitted drift
mode. Expanded by columns $h_j$ of $H$, the exact moment conditions are, for
diagonal $R = \diag(\sigma_n^2)$,
\begin{equation}
\sum_{n=1}^N h_j(t_n)\, \frac{s_n}{\sigma_n^2}\, m_n = 0,
\qquad j = 1,\ldots,p.
\tag{O4.4}\label{eq:o4-moment}
\end{equation}

\paragraph{Meaning: chopping, interleaving, randomization.}
This is the theorem behind chopping, interleaving, and randomization:
\begin{itemize}
\item \emph{chopping/paired designs} are exact when their signed effective rows
annihilate the low-order temporal moments in~\eqref{eq:o4-moment};
\item \emph{interleaving} reduces the correlation between object-carrier
directions and low-frequency gain modes;
\item \emph{random permutation} makes the cross block mean-zero under centered
exchangeable rows and gives concentration around zero.
\end{itemize}
The first two are exact only when the pattern/carrier moment equations actually
hold. ``Alternating is always optimal'' is false.

% =============================================================================
\section{Randomization corollary and flip-boundary rederivation}
\label{sec:o4-rand}
\subsection{Randomization proposition and its obstacle}
\label{sec:o4-randprop}
\begin{proposition}[Randomization concentration]
\label{prop:rand}
Under bounded centered carrier vectors and a uniformly random permutation,
standard matrix concentration bounds
\begin{equation*}
\| I_{x\beta} \|_{\mathrm{op}}
= O_{\mathbb P}\!\left( B\sqrt{ \frac{\log((K+p)/\delta)}{N} } \right),
\end{equation*}
which in turn bounds the Schur-complement loss by
\begin{equation*}
\| I_{xx} - I_{x\mid\mathrm{gain}} \|_{\mathrm{op}}
\le \frac{ \| I_{x\beta} \|_{\mathrm{op}}^2 }{ \lambda_{\min}(I_{\beta\beta}) }.
\end{equation*}
\end{proposition}

\begin{remark}[Named obstacle --- open lemma]
\label{rem:rand-obstacle}
Proposition~\ref{prop:rand} is a theorem only under explicitly bounded,
centered, permutation-sampling assumptions. It is \emph{not} yet a theorem for
arbitrary nonnegative GI patterns, because $s_n = m_n^\top x$ is scene-dependent
and can destroy centering. The named obstacle is the product $s_n m_n$: temporal
randomization randomizes a scene-dependent carrier, not a fixed zero-mean row
sequence. \emph{Closing this lemma requires the stationary-carrier hypothesis of
the companion Part~I}~\cite{companion}; that hypothesis is exactly the condition
needed to complete the proof. This remark must remain until the lemma is closed.
\end{remark}

\subsection{Relation to the existing flip boundary}
\label{sec:o4-flip}
The paired schedule can spend measurements/references to reduce $I_{x\beta}$; a
randomized schedule preserves more distinct object rows but only suppresses the
cross block statistically. The correct flip criterion is therefore equality of a
task criterion $\Phi$ applied to the two Schur complements,
\begin{equation}
\Phi\{ I_{x\mid\mathrm{gain}}^{\mathrm{pair}} \}
= \Phi\{ I_{x\mid\mathrm{gain}}^{\mathrm{rand}} \},
\tag{O4.5}\label{eq:o4-flip}
\end{equation}
under the same total-time and photon accounting. Part~I's existing $\rho^*$ flip
may be a closed-form instance of this equality; it should be rederived
from~\eqref{eq:o4-schur}, not merely re-described as a heuristic schedule
boundary. The measure-valued/KKT extension is then classical optimal
design~\cite{pukelsheim2006,dette1997,imhofwong2000} applied to atoms that carry
both a spatial pattern and a temporal slot. The novel content is the
correlated-gain information atom and the nuisance-orthogonality geometry.

% =============================================================================
\section{A model-specific temporal operating surface}
\label{sec:o2}
This section is a deliberate stub. The R26 ruling rejects a universal two-axis
``temporal ridge'' and retains only a conditional program: a decorrelation-time
ratio alone does not create an interior optimum, because gain averaging over a
frame is monotone and produces no interior extremum. An operating surface
becomes well posed only after a competing mechanism and a full resource budget
are fixed --- finite total acquisition time, read noise or switching overhead,
per-frame photon budget, pattern rank/identifiability, and the exact
within-frame gain functional. Accordingly we state the operating map as a
future, conditional target and make no ridge claim here.

\begin{conjecture}[O2 --- resource-conditioned operating surface; FUTURE / conditional]
\label{conj:o2}
Fix a specified OU-gain process, a fixed total acquisition time, and a
read-noise-plus-shot-noise GI model. Then a matched temporal operating surface
exists on a slice of the resource parameters
\begin{equation*}
r = \frac{T_f}{t_c},
\quad
\eta = \lambda t_c,
\quad
R_{\mathrm{tot}} = \frac{T_{\mathrm{total}}}{t_c},
\quad
\omega = \frac{t_{\mathrm{switch}}}{t_c},
\quad
N/K,
\end{equation*}
together with a read-noise-to-shot-noise ratio. Any ridge, its constant, and
even its existence are model-dependent and must be derived only after these
preconditions are fixed. The named obstacle is the coupled growth/decline of
pattern rank, gain posterior uncertainty, and per-frame photon information:
there is no scalar renewal reduction analogous to dead time. This is a
conjecture, not a theorem, and it must not enter the title or the abstract until
a concrete resource-conditioned law is derived.
\end{conjecture}

% =============================================================================
\section{Reach: gain drift plus dead-time counting}
\label{sec:o5}
Let the complete latent record contain the gain path, photon arrivals, and
detector state. Conditional-score operators compose through nested observation
channels by the tower property, so the controlled-score Gramian architecture
survives. What does \emph{not} survive is a scalar product law. If a frame's
incident rate is $a_n s_n$ and the counter response is nonlinear, the score
depends on
\begin{equation*}
\frac{\partial}{\partial x} \log p_{\mathrm{count}}(Y_n\mid a_n s_n),
\end{equation*}
so the posterior uncertainty of $a_n$ is weighted by the derivative and
curvature of the dead-time count law. High-gain states are preferentially
compressed or saturated. Consequently:
\begin{itemize}
\item gain uncertainty and dead-time loss are coupled;
\item temporal carriers are clipped most strongly exactly when $a_n$ is large;
\item a correction trained for an analog linear bucket may become biased under
photon-counting saturation;
\item the best temporal schedule may shift because extreme-gain frames carry
less marginal object information.
\end{itemize}
The exact O1 identity~\eqref{eq:o1-fisher} still holds if the complete score
includes both hidden mechanisms. The first composite treatment should derive one
explicit Poisson-renewal specialization; it must \textbf{not} claim a
mechanism-by-mechanism multiplicative efficiency law.

% =============================================================================
\section*{Code and data availability}
Numerical checks of the O4.2 Schur formula and the O4.4 orthogonality
conditions, and any released code and evidence bundles, are available
at~\SPH{repo URL wording --- user decision}.

\section*{Acknowledgments}
\SPH{funding/acknowledgments --- user decision}

\nocite{*}
\bibliographystyle{unsrt}
\bibliography{refs}

\end{document}
